\subsection{Evidence eligibility and claim boundary}
The empirical evidence in this paper consists only of the completed E7 codec
benchmark and E8 cold-start replay under the frozen
\texttt{strict\_evaluation\_v1} contract. These are separate audits with
different data sources and tasks; they are not pooled as repeated observations
of one physical mechanism. Historical exploratory comparisons and invalidated
mechanism-chain artifacts are excluded from configuration selection, figures,
tables, and conclusions. In particular, a protocol or runner invalidation is
an integrity event rather than a negative measurement of codec performance.

Accordingly, the paper makes claims only about the declared post-compensation
software representations, data splits, codec families, and replay features.
It does not estimate population probability of detection, calibrated field
false-alarm rates, cross-structure performance, hardware energy, embedded
memory, throughput, real-time latency, or deployment readiness.

For the reviewed submission snapshot, a source manifest hashes every TeX
input, bibliography, and included result figure together with the frozen
protocol and result artifacts. A separate build receipt records one rendered
PDF against that committed source snapshot. The accompanying read-only checks
reject source drift and direct historical identifiers in the manuscript; they
identify the reviewed files, not the original execution chronology or any
broader scientific claim.

\subsection{Scope and shared representation}
The system boundary is post-compensation waveform coding. Each method receives
the same 40~kHz, high-pass, unit-energy-normalized OBS+BSS residual record
and reconstructs it in software. The benchmark does not measure original ADC
acquisition, MCU execution time, energy, latency, packet loss, or embedded
memory. A signed-16-bit quantizer is fitted from the maximum absolute healthy
training residual; its scale is then fixed for validation, test, D04, and
D24 records.

\subsection{Hard-cap codecs}
Each monitoring record contains 66 paths. The payload metric is the actual
serialized byte length of the full record packet, not a count of retained
samples. We compare four codecs.

\paragraph{Bounded SoD.}
For each path, SoD writes an initial quantized level followed by variable
length timestamp gaps and signed level changes. A target-dependent path
budget includes the path framing. Encoding stops only before an event that
would exceed that budget, and the decoder holds the final transmitted level
for the remaining samples. Thus every bounded packet remains decodable and
has a proven record-level upper bound.

\paragraph{Uniform linear interpolation, PCA, and Haar DWT.}
The uniform codec stores a fixed-stride sequence and reconstructs it by
linear interpolation. PCA uses a basis fitted from healthy training records
only, while Haar DWT keeps a fixed number of transform coefficients. For
these three codecs, only candidates with a proven packet bound no greater
than the target are eligible. Decoder model bytes, including the PCA basis,
are reported separately from the per-record payload rather than silently
treated as free.

\subsection{Codec selection and scoring}
The protocol fixes capacities of 2,048, 4,096, 8,192, and 16,384 bytes per
record. For fixed-packet codecs, the eligible configuration closest to the
target bound is selected. For bounded SoD, healthy-validation mean payload
closest to the target is selected; a payload tie is resolved by
code-domain reconstruction MSE on eight chronologically spaced healthy
validation records, followed by the frozen candidate order. No damage labels
or test records enter this selection.

Each decoded record is summarized by residual reconstruction energy aggregated
over all paths. We compute held-out record ROC AUC and a stratified bootstrap
95\% interval. As a temperature-control diagnostic, healthy and damage
records are matched one-to-one by temperature using a Hungarian assignment,
with no record reuse and a maximum permitted difference of
0.25~$^\circ$C. Matching is reported alongside, rather than substituted for,
the full record-level AUC.

\subsection{Cold-start alarm replay}
The alarm evaluation uses an independent software replay. A reference bank,
per-path robust normalization, Level-A SoD deltas, and the complete
nine-value threshold grid are fitted only from healthy March 2021 data.
Dense residual energy and Level-A SoD event count are then scored over all
April 2021 records before labels are read for reporting.

An alarm incident is a run of threshold exceedances with gaps no larger than
30 minutes. For every frozen threshold, we report false calls per day on the
labelled healthy April prefix, the delay to the first newly started
post-onset incident, and post-onset record and day exceedance coverage. An
incident active before the onset is flagged and cannot be credited as a new
detection. The single observed labelled transition is therefore not called a
probability of detection.
